\begin{table}[!htbp]\small\setlength{\tabcolsep}{4pt}
\caption{Composantes de la réforme foncière (Md€ de 2025)}\label{tab:t3}
\begin{tabularx}{\linewidth}{@{}Xrrc@{}}
\toprule
\textbf{Composante} & \textbf{2032} & \textbf{À terme} & \textbf{Statut} \\
\midrule
Taxe nationale sur la valeur du sol (taux de 0,24 %) & +15,0 & +15,0 & HYP \\
Déductibilité de la taxe sur le sol, calendrier de recouvrement & $-$1,8 & $-$1,5 & DED \\
Droits de mutation divisés par deux & $-$8,1 & $-$8,1 & OBS \\
Taxe communale sur la valeur de marché, en deux parts & 0,0 & 0,0 & HYP \\
\midrule \textbf{Solde public} & \textbf{+5,2} & \textbf{+5,4} & \\
\bottomrule
\end{tabularx}
\par\vspace{2pt}{\raggedright\scriptsize\color{gris}Source : valeurs foncières des ventes (DVF), comptes de patrimoine de l'Insee ; droits de mutation au taux permanent, 16,1 Md€ en 2025 (OFGL, DGFiP) ; calculs des auteurs. Assiette du sol : 6\,156 Md€ après la baisse de valeur que provoque la taxe (construite, non observée). La taxe communale garde sa recette de 2025.\par}
\end{table}
